\section{Discussion}\label{sec:discussion}

\subsection{A shared functional response can conceal different mechanisms}\label{sec:discussion-mechanisms}

The central empirical result is a failed mechanism-transfer hypothesis. GPT-2 IOI L10H2 becomes more useful after removal of the canonical primary Name Movers and shows target-positive/rival-negative score contributions. Pythia-160M induction L8H2 also becomes more useful after its primary intervention, yet the fine score signs reverse. The two settings therefore share the functional label \emph{conditional recruitment} without sharing the tested fine-grained mechanism.

This distinction matters because recruitment is defined by a change in causal usefulness, not by a particular internal adjustment. GPT-2 L10H10 makes the same point within a single model and task: it increases target attention and exhibits the same target-positive/rival-negative score organization as recruited L10H2, yet its overall conditional gain is negative. Neither an attention increase nor the sign of one FALI channel is sufficient to define recruitment. The output-grounded candidate effect defines the functional response; FALI then allocates the cross-state change under the specified intervention game.

The cross-setting reversal should not be over-interpreted. Model family, task, head pair, perturbation, and metric all differ between GPT-2 IOI and Pythia induction. The experiment therefore shows that the GPT-2 fine-score pattern is not a universal consequence of positive conditional recruitment in the tested settings. It does not identify which of those differences caused the reversal, nor does it rule out deeper common structure.

\subsection{Where FALI fits in the interpretability workflow}\label{sec:discussion-positioning}

FALI is best viewed as a mechanism-resolution stage downstream of localization. Activation and path patching identify behaviorally relevant activations or paths \citep{goldowskydill2023pathpatching,zhang2024activationpatching,heimersheim2024activationpatching}; conditional methods such as CoAx identify components whose importance changes after a primary intervention \citep{gong2026coax}; interaction-aware mediation analyses establish that component effects can depend on the surrounding network \citep{vaidyanathan2026curse}. FALI begins once a selected attention head and target edge are already specified, and asks which score, value, or surrounding-context changes account for its altered usefulness.

The framework also does not claim novelty for QK/OV separation or Owen attribution. Routing-versus-content distinctions are established in transformer-circuit analysis and more recent causal decompositions \citep{elhage2021mathematical,franco2025attentioncausal,wang2026fewshot}. Standard Owen values provide interaction-aware accounting under a fixed coalition structure. FALI's contribution is to apply these ingredients to a different estimand: the intervention-induced change in one selected head's output-grounded causal usefulness, with target/rival/other resolution inside SCORE and an explicit surrounding-context factor.

This framing suggests a practical pipeline: first localize a component or detect conditional importance; then, when the scientific question is \emph{why that selected attention head's usefulness changed}, apply FALI under a prospectively specified intervention game. The output should be treated as a discriminative counterfactual explanation, not as an automatically discovered natural pathway.
